LLM agents are proliferating across enterprise and open ecosystems, each advertising tools, APIs, and knowledge-domain capabilities. Yet no lightweight, verifiable, lifecycle-aware registry exists to govern what a capability claims, how it evolves, and whether it remains trustworthy. Existing governance layers---KYA for permissions and AgentSafe for architecture-level oversight---leave this gap unaddressed. ADL Lite fills it as an event-first capability-lifecycle registry. Drawing on Wittgenstein's thesis that ``the world is the totality of facts, not of things,'' ADL Lite treats events---not objects---as the fundamental ontological primitive, while distinguishing occurrents (events as they occur) from their immutable records (EventChains). Capabilities are modeled as append-only, cryptographically linked EventChains, with all lifecycle state (status, confidence, validators) derived exclusively from event history through deterministic functions. Every capability claim, validation, deprecation, and fork is an auditable event; this design eliminates stored mutable state and the consistency defects it introduces.

We provide a two-level ontological analysis that situates ADL Lite within foundational ontologies (BFO, DOLCE, UFO), formalizes identity conditions and three forms of ontological dependence, and defines a decidable precondition language for lifecycle transitions. We further provide a complete formal semantics, together with seven proved properties including determinism, confluence under fork, and CRDT convergence.

We empirically validated the registry on 201 EventChains derived from the IBM Anti-Money Laundering dataset (9,300 events), confirming zero integrity failures and linear scalability under architectural stress tests. These findings are supplemented by concurrency stress tests, adversarial integrity trials, and comparative baselines against nanopublications and Git-native systems. This paper is positioned as a framework paper establishing the ontological and formal foundations of capability-lifecycle governance. The registry is implemented as a pip-installable Python toolkit operating on Markdown files in standard Git repositories, targeting lightweight deployment scenarios where full OWL/SPARQL infrastructure is unavailable or unnecessary.
